% id: 134-13
% book: 베이직쎈 미적분1 · 08 정적분
% section: 유형 085 정적분으로 정의된 함수의 극대·극소
% figure: tikz:fig-134-13
% answer: 
% answer_source: 
% variant_level: 0 (원본 전사)
% 이 파일은 build.mjs 가 items.json 에서 만든다. 고칠 때는 items.json 을 고친다.
\smash{\raisebox{4.5mm}{\dmkichul{베이직쎈 미적분1 134쪽 13번}}}
\noindent\begin{minipage}[t]{0.58\linewidth}\vspace{0pt}\raggedright 이차함수 $y=f(x)$의 그래프가 오른쪽 그림과 같을 때, $F(x)=\displaystyle\int_{0}^{x} f(t)\mkern3mu dt$를 만족시키는 함수 $F(x)$에 대하여 다음을 구하시오.\end{minipage}\hfill\begin{minipage}[t]{0.4\linewidth}\vspace{0pt}\centering \resizebox{\ifdim\width>\linewidth\linewidth\else\width\fi}{!}{% 134-13(134쪽) — 이차함수 $y=f(x)$ 의 그래프($x$축과 $x=-2$, $x=0$ 에서 만나고 꼭짓점이 $(-1,\,-1)$ 인 아래로 볼록한 포물선). 스캔 화질이 낮아 TikZ 로 다시 그림.
% 원문에 있는 것만 옮긴다: 곡선 · 꼭짓점에서 $x$축·$y$축으로 가는 안내 점선 두 줄 · 라벨 $-2$, $-1$(축 위), $-1$($y$값), O, x, y, y=f(x). 원문에 없는 눈금·값은 적지 않는다.
\begin{tikzpicture}[x=0.5cm, y=0.5cm, line width=0.6pt]
  \draw[->, >=stealth, line width=0.7pt] (-3.15,0) -- (1.5,0) node[below=1pt, font=\small] {$x$};
  \draw[->, >=stealth, line width=0.7pt] (0,-1.75) -- (0,2.6) node[left=1pt, font=\small] {$y$};
  \draw[densely dashed, line width=0.4pt] (-1,-1) -- (-1,0);
  \draw[densely dashed, line width=0.4pt] (-1,-1) -- (0,-1);
  \draw[domain=-2.55:0.55, samples=80, smooth] plot (\x, {(\x)*(\x)+2*(\x)});
  \node[below=1pt, font=\small] at (-2.25,0) {$-2$};
  \node[above=1pt, font=\small] at (-1,0) {$-1$};
  \node[below right=-1pt and 0pt, font=\small] at (0,0) {$\pt{O}$};
  \node[right=2pt, font=\small] at (0,-1) {$-1$};
  \node[anchor=west, inner sep=1pt, font=\small] at (0.2,2.4) {$y=f(x)$};
\end{tikzpicture}
}\end{minipage}\par
\dmsubprob{1} 함수 $f(x)$
\dmsubprob{2} $F(x)$가 극댓값을 갖는 $x$의 값
\dmsubprob{3} $F(x)$의 극댓값
